\chapter{LTI 系の基礎}


\section{システムとして信号処理を見る}

これまでの章では，主に信号そのものを見てきました．
本章では，信号を入力すると別の信号を出力する仕組みを\textbf{システム}として考えます．

システムを
\[
y[n] = \mathcal{T}\{x[n]\}
\]
と書くことにします．
ここで $\mathcal{T}$ は変換規則そのものを表します．

例えば，
\begin{itemize}
\item ノイズを減らす平滑化フィルタ
\item 画像のぼかしやエッジ抽出
\item 音声の残響付与
\item ニューラルネットワークの 1 層
\end{itemize}
は，いずれも広い意味で入力信号を出力信号へ写すシステムです．

その中でも特に重要なのが，\textbf{線形時不変 (Linear Time-Invariant; LTI) 系}です．
LTI 系は解析しやすく，多くのフィルタや信号処理アルゴリズムの基礎になります．


\section{線形性}

システム $\mathcal{T}$ が\textbf{線形}であるとは，任意の信号 $x_1[n], x_2[n]$ と定数 $a, b$ に対して
\[
\mathcal{T}\{a x_1[n] + b x_2[n]\}
= a\mathcal{T}\{x_1[n]\} + b\mathcal{T}\{x_2[n]\}
\]
が成り立つことです．

これは
\begin{itemize}
\item 入力を足したとき，出力も足し合わせで表せる
\item 入力を定数倍したとき，出力も同じ定数倍になる
\end{itemize}
という 2 つの性質をまとめたものです．

線形性があると，複雑な入力を簡単な成分に分けて考えられます．
スペクトル解析や畳み込みが有効なのは，この分解可能性のおかげです．


\section{時不変性}

システム $\mathcal{T}$ が\textbf{時不変}であるとは，入力を $n_0$ サンプル遅らせたとき，
出力も同じだけ遅れることです．
式で書くと，
\[
\mathcal{T}\{x[n-n_0]\} = y[n-n_0]
\]
です．

この性質は，「いつ入力したか」ではなく「どのような形の入力か」にだけシステムが反応することを意味します．
同じ拍手音でも，1 秒後に入れても 2 秒後に入れても，応答波形の形は変わらず位置だけがずれる，というのが時不変の直感です．


\section{インパルス応答}

LTI 系に単位インパルス $\delta[n]$ を入力したときの出力を
\[
h[n] = \mathcal{T}\{\delta[n]\}
\]
と定義し，これを\textbf{インパルス応答}と呼びます．

インパルス応答が重要なのは，LTI 系では任意の入力に対する出力を $h[n]$ だけで決められるからです．
これは，任意の離散信号が
\[
x[n] = \sum_{m=-\infty}^{\infty} x[m]\delta[n-m]
\]
と表せることから分かります．

右辺は「各時刻 $m$ に置かれたシフト付きインパルスを，重み $x[m]$ で足し合わせたもの」です．
LTI 系にこの表現を入れると，線形性と時不変性によって
\[
y[n] = \sum_{m=-\infty}^{\infty} x[m] h[n-m]
\]
が得られます．

これが\textbf{畳み込み和}です．


\section{畳み込み}

離散時間の畳み込みを
\[
y[n] = (x * h)[n] = \sum_{m=-\infty}^{\infty} x[m] h[n-m]
\]
と書きます．

畳み込みの計算では，$h[n]$ を反転して平行移動し，入力 $x[m]$ と重ね合わせながら総和を取る，という見方が有用です．
公式だけ暗記するのではなく，
\begin{itemize}
\item 反転
\item シフト
\item 各点での積
\item 総和
\end{itemize}
の流れを頭の中で再現できるようにしておく必要があります．

有限長の FIR フィルタなら，実際には和の範囲は有限個の項だけになります．
例えば
\[
h[0] = \frac{1}{3}, \quad h[1] = \frac{1}{3}, \quad h[2] = \frac{1}{3}
\]
だけが非ゼロなら，
\[
y[n] = \frac{x[n] + x[n-1] + x[n-2]}{3}
\]
となり，3 点移動平均フィルタになります．


\section{差分方程式による表現}

LTI 系は，差分方程式で書かれることも多いです．
例えば
\[
y[n] = x[n] + 0.5x[n-1]
\]
は FIR 系です．

一方，
\[
y[n] = 0.8 y[n-1] + x[n]
\]
のように過去の出力を使う系は IIR 系です．

FIR は有限長のインパルス応答を持ち，実装や安定性の扱いが比較的簡単です．
IIR は少ない係数で鋭い周波数特性を作りやすい一方で，安定性や初期条件の扱いが重要になります．


\section{因果性と BIBO 安定性}

\subsection{因果性}

実時間で動作するシステムでは，現在の出力が未来の入力に依存してはいけません．
この性質を\textbf{因果性}と呼びます．

LTI 系が因果的であるためには，
\[
h[n] = 0 \quad (n < 0)
\]
であることが必要です．
つまり，インパルスを入れる前に応答が始まってはいけません．


\subsection{BIBO 安定性}

有界な入力に対して常に有界な出力が得られる性質を\textbf{BIBO 安定}と呼びます．
離散時間 LTI 系では
\[
\sum_{n=-\infty}^{\infty} |h[n]| < \infty
\]
であれば BIBO 安定です．

直感的には，インパルス応答の総量が無限に大きくならないなら，入力を有限範囲に抑えたとき出力も暴走しない，ということです．


\section{周波数応答}

複素指数信号
\[
x[n] = \exp(\mathrm{j}\omega n)
\]
を LTI 系に入れると，出力は同じ周波数のまま定数倍されます．
すなわち
\[
y[n] = H(\mathrm{e}^{\mathrm{j}\omega}) \exp(\mathrm{j}\omega n)
\]
と書けます．

ここで
\[
H(\mathrm{e}^{\mathrm{j}\omega}) = \sum_{n=-\infty}^{\infty} h[n]\exp(-\mathrm{j}\omega n)
\]
を\textbf{周波数応答}と呼びます．

$|H(\mathrm{e}^{\mathrm{j}\omega})|$ は周波数ごとの増幅率，
$\angle H(\mathrm{e}^{\mathrm{j}\omega})$ は位相のずれを表します．
ローパスフィルタ，高域強調，帯域通過といった設計は，この周波数応答を狙って行います．


\section{Python で移動平均と IIR を比べる}

以下のコードでは，ノイズを含む信号に対して，FIR の移動平均と 1 次 IIR を適用します．

\begin{lstlisting}[style=codecell,language=Python]
from pathlib import Path

import matplotlib.pyplot as plt
import numpy as np

output_dir = Path("outputs/ch15")
output_dir.mkdir(parents=True, exist_ok=True)

rng = np.random.default_rng(0)
n = np.arange(200)
x = np.sin(2 * np.pi * 0.03 * n) + 0.4 * rng.standard_normal(size=n.shape)

h = np.ones(5) / 5
y_fir = np.convolve(x, h, mode="same")

y_iir = np.zeros_like(x)
alpha = 0.8
for i in range(len(x)):
    prev = y_iir[i - 1] if i > 0 else 0.0
    y_iir[i] = alpha * prev + (1 - alpha) * x[i]

plt.figure(figsize=(8, 3))
plt.plot(n, x, label="input", alpha=0.5)
plt.plot(n, y_fir, label="FIR moving average", linewidth=2)
plt.plot(n, y_iir, label="first-order IIR", linewidth=2)
plt.xlabel("sample index n")
plt.ylabel("amplitude")
plt.legend()
plt.tight_layout()
plt.savefig(output_dir / "fir_vs_iir.png", dpi=150)
plt.close()
\end{lstlisting}

この例では，どちらも平滑化を行いますが，性質は同じではありません．
\begin{itemize}
\item FIR 移動平均は有限個の過去サンプルだけを見る
\item IIR は過去の出力を再帰的に使うので，より長い記憶を持つ
\item 同じ「なめらかにする」処理でも，遅延や立ち上がりの鋭さが異なる
\end{itemize}

この差は，周波数応答を見るとさらに明確になります．


\section{この章で押さえるべき点}

\begin{itemize}
\item LTI 系は線形性と時不変性を持つシステムである
\item LTI 系はインパルス応答 $h[n]$ だけで記述できる
\item 任意入力への出力は畳み込み和 $y[n] = (x*h)[n]$ で与えられる
\item 因果性は $h[n]=0$ for $n<0$，BIBO 安定性は $\sum_n |h[n]| < \infty$ と関係する
\item 周波数応答 $H(\mathrm{e}^{\mathrm{j}\omega})$ を見ると，フィルタが各周波数をどう変えるか分かる
\end{itemize}

本章まで理解できれば，離散信号，スペクトル，LTI 系という信号処理の中核概念が一通りつながります．
この基盤の上に，STFT，画像フィルタ，音響特徴量抽出，畳み込みニューラルネットワークの理解を積み上げていけます．
